% Part: history
% Chapter: biographies
% Section: ernst-zermelo

\documentclass[../../../include/open-logic-section]{subfiles}

\begin{document}

\olfileid{his}{bio}{zer}

\olsection{ارنست زرملو}

\olphoto{zermelo-ernst}{Ernst Zermelo}

ارنست زرملو در ۲۷ ژوئیهٔ ۱۸۷۱ در برلینِ آلمان زاده شد. او پنج خواهر
داشت، اما خانواده‌اش از بیماری رنج می‌بردند و تنها سه تن از خواهرانش به
بزرگسالی رسیدند. پدر و مادرش نیز هنگامی که او جوان بود درگذشتند و زرملو و
خواهران و برادرانش در هفده‌سالگی او یتیم شدند. زرملو علاقه‌ای عمیق به هنر
و به‌ویژه شعر داشت. او به تیزهوشی، بذله‌گویی و نگاه انتقادی شهره بود.
مشهورترین دستاوردهای ریاضی‌اش معرفی اصل انتخاب در ۱۹۰۴ و اصل‌موضوعی‌کردن
نظریهٔ مجموعه‌ها در ۱۹۰۸ است.

علایق دانشگاهی زرملو متنوع بود. او درس‌هایی در فیزیک، ریاضیات و فلسفه
گذراند. زرملو در سال ۱۸۹۴، زیر نظر هرمان شوارتس، رسالهٔ
\emph{پژوهش‌هایی در حساب تغییرات} را در دانشگاه برلین به پایان رساند. در
سال ۱۸۹۷ تصمیم گرفت تحصیلاتش را در دانشگاه گوتینگن ادامه دهد؛ جایی که
کارهای بنیانی دیوید هیلبرت تأثیر بسیاری بر او گذاشت. در سال ۱۸۹۹ صلاحیت
احراز کرسی استادی را به دست آورد، اما تا یازده سال بعد کرسی‌ای نگرفت؛
شاید به سبب رفتار نامتعارف و «شتاب عصبی» او.

زرملو سرانجام در سال ۱۹۱۰ کرسی استادیِ باحقوقی در دانشگاه زوریخ گرفت،
اما ابتلا به سل او را در ۱۹۱۶ به بازنشستگی واداشت. پس از بهبودی، در سال
۱۹۲۱ کرسی افتخاری دانشگاه فرایبورگ به او داده شد. در این دوره بر مبانی
ریاضیات کار کرد. نوشته‌های تورالف اسکولم و کورت گودل او را آزرده می‌کرد و
در مقالاتش آشکارا از رویکردهای آنان انتقاد کرد. در سال ۱۹۳۵، به سبب
نامحبوب‌بودن و مخالفتش با به‌قدرت‌رسیدن هیتلر در آلمان، از مقام خود در
فرایبورگ برکنار شد.

سال‌های پایانی زندگی زرملو با انزوا همراه بود. پس از برکناری در ۱۹۳۵،
ریاضیات را کنار گذاشت و به روستا رفت و ساده زیست. در ۱۹۴۴ ازدواج کرد و
چون بینایی‌اش رو به زوال بود کاملاً به همسرش وابسته شد. زرملو تا سال ۱۹۵۱
به‌کلی نابینا شد. او در ۲۱ مه ۱۹۵۳ در گونترستالِ آلمان درگذشت.

\begin{reading}
برای زندگی‌نامهٔ کامل زرملو، بنگرید به \citet{Ebbinghaus2015}. مقالات
بنیادین او در سال‌های ۱۹۰۴ و ۱۹۰۸ به زبان اصلی آلمانی در
\citep{Zermelo1904,Zermelo1908} در دسترس‌اند. مجموعه‌آثار زرملو، از جمله
نوشته‌هایش دربارهٔ فیزیک، با ترجمهٔ انگلیسی در
\citep{Ebbinghaus2010,Ebbinghaus2013} منتشر شده‌اند.
\end{reading}

\end{document}
